\section{Signed gaps and paired partitions (E1)}\label{sec:signedgap}

The constructions in this section use actual integer colour degrees.
In particular, a nonzero difference of four is useful, although it is
invisible modulo four. No assumption about an unknown initial colouring
is made. The results of Sections~\ref{sec:signedgap}--\ref{sec:rootports}
are independent of the subcubic proof.

\begin{lemma}[A signed gap with at most two occurrences]\label{lem:signedgap}
Let $L_i(y)=\sum_{e\in E_i}a_{ie}y_e$, where $a_{ie}\in\{-1,1\}$,
every $E_i$ contains at least four distinct variables, and every variable
belongs to at most two different sets $E_i$. There is a choice
$y_e\in\{-1,1\}$ for which $|L_i(y)|\ge2$ for every $i$.
Such a choice can be found in polynomial time.
\end{lemma}
\begin{proof}
Make an auxiliary vertex for each form. A variable occurring twice is an
edge between its two forms, with label $\tau_e=a_{ie}a_{je}$; parallel
edges are allowed, but there are no loops. A variable occurring once is
a boundary variable at its form.

Choose signs $\sigma_i\in\{-1,1\}$ and call an auxiliary edge good when
$\tau_e=\sigma_i\sigma_j$. Start with any signs and flip a vertex whenever
it has more bad than good internal edges. Each flip strictly increases
the total number of good edges, so there are at most as many flips as
internal edges. At termination, the number of good internal edges at
each vertex is at least its number $b_i$ of bad edges. Count boundary
variables as good and write $g_i$ for this enlarged good count. Then
\[
 m_i=g_i-b_i=|E_i|-2b_i\ge0.
\]
For a boundary variable put $y_e=\sigma_i a_{ie}$; for a good edge put
$y_e=\sigma_i a_{ie}=\sigma_j a_{je}$. A bad edge contributes opposite
values at its ends after multiplying the two rows by $\sigma_i$ and
$\sigma_j$. We may therefore select which end receives $+1$. If $o_i$
bad edges select $i$ as their positive end, then
\begin{equation}\label{eq:signedmargin}
 \sigma_i L_i=g_i-b_i+2o_i=m_i+2o_i.
\end{equation}
Call a vertex critical if $m_i\in\{0,1\}$. Since $|E_i|\ge4$, every
critical vertex has bad-edge degree at least two. In a component of the
bad-edge graph with a noncritical vertex, root a spanning tree at such
a vertex and select each nonroot vertex as the positive end of its
parent edge. If all vertices of a component are critical, its minimum
degree is at least two, so it contains a cycle, possibly a pair of
parallel edges. Select positive ends cyclically around that cycle and
use attached rooted trees to give every remaining vertex a positive
parent edge. Set the positive ends of unused edges arbitrarily.
Every critical vertex now has $o_i\ge1$, and every noncritical vertex
has $m_i\ge2$. Equation~\eqref{eq:signedmargin} proves the claim.
The local flips and all tree and cycle operations have polynomial bounds.
\end{proof}

Odd-length rows are permitted; their sums have odd parity, so their
absolute values are actually at least three. This distinction will
remove a boundary obstruction on an even-order path.

\begin{theorem}[Paired vertex partition]\label{thm:integerpairs}
Let $V(G)=S\sqcup T$ be a supplied partition. Suppose that both induced
graphs have maximum degree at most one, that $G[T]$ has no isolated
vertex, and that every edge $uv$ internal to either side satisfies
\[
 \min(d(u),d(v))\ge3,\qquad |d(u)-d(v)|\le1.
\]
Then $\lir(\doub G)\le2$. The mixed original edges can be exactly
$E(G[T])$. No degree bound is imposed on an internally isolated vertex
of $S$, and there is no upper bound on the maximum degree of $G$.
\end{theorem}
\begin{proof}
Give edges of $G[S]$ weight zero, edges of $G[T]$ weight one, and each
crossing edge $e$ weight $1+y_e$, where $y_e\in\{-1,1\}$. For each
internal pair $uv$ form
\[
 L_{uv}=\sum_{\substack{e\text{ crossing}\\e\ni u}}y_e
       -\sum_{\substack{e\text{ crossing}\\e\ni v}}y_e.
\]
The variables in a row are distinct and number $d(u)+d(v)-2\ge4$.
Each crossing edge occurs in at most one row at each end, hence at most
twice altogether. Apply Lemma~\ref{lem:signedgap}. The equal internal
contributions at $u,v$ give
\begin{equation}\label{eq:integerpairdifference}
 R(u)-R(v)=d(u)-d(v)+L_{uv},\qquad
 B(u)-B(v)=d(u)-d(v)-L_{uv}.
\end{equation}
Both differences are nonzero. On crossing edges, vertices of $S$ have
even red and blue degrees, whereas vertices of $T$ have odd red and
blue degrees. Thus all crossing edges are legal as well.
\end{proof}

The assumption on degree differences concerns the endpoints of the
internal pairs, not the degree range of the whole graph. The artifact
includes a connected nonbipartite example with a leaf and degrees
ranging from one to six.

\subsection{An integer repair with a fixed mixed set}

To see why retaining integer values matters, take the cube on
$0,\ldots,7$, with adjacency given by exclusive-or differences $1,2,4$,
and add $03$. Put $S=\{0,1,6,7\}$, $T=\{2,3,4,5\}$ and
$M=\{23,45\}$. The following weights and degrees are a complete witness.
\begin{center}\small
\begin{tabular}{c|rrrrrrrrrrrrr}
$e$&01&02&04&13&15&23&26&37&45&46&57&67&03\\\hline
$x_e$&0&2&2&0&0&1&2&0&1&2&0&0&0
\end{tabular}\par\medskip
\begin{tabular}{c|rrrrrrrr}
$v$&0&1&2&3&4&5&6&7\\\hline
$R(v)$&4&0&5&1&5&1&4&0\\
$B(v)$&4&6&1&7&1&5&2&6
\end{tabular}
\end{center}
With this same $M$, the stronger demand that both differences on every
mixed edge be nonzero modulo four is impossible: on $23$ the red degrees
are odd, while the ordinary degrees differ by one. A nonzero red
difference modulo four would have to be two, making the blue difference
zero modulo four. In the displayed legal colouring the actual red
difference is four and the blue difference is minus six.

There is also an infinite global obstruction to that stronger demand.
For $t\ge1$, take $H_t=C_{4t}$ with its antipodal matching. Subdivide
every edge once, let $S$ be the $4t$ old vertices and $T$ the $6t$ edge
vertices, and add a perfect matching inside each side. The resulting
graph $F_t$ has $10t$ vertices and $17t$ edges; its $S$ degrees are four
and its $T$ degrees are three. It is connected and nonbipartite: an
internal $S$ matching edge closes an even path in the subdivision into
an odd cycle. Theorem~\ref{thm:integerpairs} colours every $F_t$ with
exactly the $T$ matching mixed.

For the same mixed set, the stronger modulo-four condition is feasible
if and only if $t$ is even. Indeed, write
$R(v)\equiv p_v+2q_v\pmod4$, with $p=0$ on $S$ and $p=1$ on $T$.
Every internal pair must have opposite $q$ values. The $5t$ pairs force
$\sum_vq_v\equiv5t\pmod2$. On the other hand, writing every nonmixed
weight as $2z_e$ gives $q_v\equiv\sum_{e\ni v,\,e\notin M}z_e\pmod2$,
so the same sum is even. Conversely, when $t$ is even, assign opposite
$q$ values on each pair. Their even total can be realised as vertex
parities by a spanning-tree leaf elimination in the connected graph
$F_t-M$. The resulting weights satisfy the stronger condition.
Thus every odd $t$ supplies an integer repair on the same graph and
with the same mixed set, rather than a repair obtained by changing
the instance or the interface.

\section{Path and cycle sides in degree three and four (E2)}\label{sec:pathcyclepart}

Here and in Section~\ref{sec:isolatepart}, every ordinary degree is three
or four. An internal degree always means degree in an induced side of
the specified partition.

\begin{lemma}[A useful maximum-cut shape]\label{lem:biasedcut}
Every connected graph with $3\le d(v)\le4$ has a partition $S\sqcup T$
whose crossing graph is connected, such that $G[S]$ has maximum degree
two, every internally degree-two vertex has ordinary degree four,
and $G[T]$ is a matching together with isolated vertices.
\end{lemma}
\begin{proof}
Among maximum cuts choose one maximising the number of ordinary
degree-four vertices in $S$. A vertex cannot have more than half of
its neighbours on its own side. Thus ordinary degree-three vertices
have internal degree at most one and degree-four vertices at most two.
If a degree-four vertex in $T$ had internal degree two, flipping it
would preserve the cut size and improve the secondary objective.
Consequently $G[T]$ has maximum degree one.
If the crossing graph were disconnected, some original same-side edge
would join two of its components, since $G$ is connected. Flipping all
vertices of one such component loses no crossing edge and gains all its
boundary edges, contrary to maximum cut size.
\end{proof}

This existence argument does not supply a polynomial-time maximum-cut
algorithm. The colouring theorems below assume their partitions are
given. They do not prove that the isolated vertices permitted by
Lemma~\ref{lem:biasedcut} can always be handled.

\begin{theorem}[Path/cycle side against a matching]\label{thm:pathcyclepart}
Suppose $d(v)\in\{3,4\}$ and a supplied partition $S\sqcup T$ satisfies:
$G[T]$ is a matching covering $T$, $\Delta(G[S])\le2$, and every
internally degree-two vertex of $S$ has ordinary degree four.
Then $\lir(\doub G)\le2$, with mixed set exactly $E(G[T])$.
The colouring can be constructed in polynomial time.
\end{theorem}

We prove the local completion facts first. Give every crossing edge
weight $1+y_e$ with $y_e\in\{-1,1\}$, every $T$ matching edge weight one,
and each $S$ internal edge weight $2z_e$ with $z_e\in\{0,1\}$.
The crossing comparisons are automatically separated by parity.
For every $T$ pair impose the row condition $|L_{uv}|\ge2$ from
\eqref{eq:integerpairdifference}. It suffices to add a matching of
selected $S$ edges to these rows, then complete the internal pure edges.

\subsection{Binary relations}

A relation on two bits has full projections if either variable, when
fixed, has some compatible value for the other. Equality, inequality,
the exclusion of one ordered pair, and the unrestricted relation all
have this property. On a variable path, full projections permit an
extension from any one prescribed bit. On a variable cycle consisting
only of equality and inequality relations, consistency is equivalent
to an even number of inequalities. A variable cycle of full-projection
relations is also consistent if at least one relation allows three or
four pairs. To prove the last claim, compose the relations around the
cycle. Full projections are preserved under composition; composing a
relation with at least three pairs with any full-projection relation
still gives at least three pairs. The resulting relation on the first
variable contains a diagonal pair. These statements include a two-edge
cycle of parallel relations.

\subsection{Completing a cycle}

Let an $S$ cycle be $v_0\ldots v_{n-1}v_0$ and put $e_i=v_iv_{i+1}$,
with indices modulo $n$. Each vertex has two crossing edges. Let $s_i$
be its number of wholly red crossing edges. Then
\[
 R(v_i)/2=s_i+z_{i-1}+z_i,
 \qquad B(v_i)/2=4-s_i-z_{i-1}-z_i.
\]
Whatever the colour of $e_i$, its legality is exactly
\begin{equation}\label{eq:cyclebits}
 s_i-s_{i+1}+z_{i-1}-z_{i+1}\ne0.
\end{equation}
When the first difference is zero this is inequality of the two bits;
when it is $1$ or $-1$ it excludes one ordered pair; when its absolute
value is at least two it imposes no restriction. The variable graph
uses step two around the cycle. For odd $n$ it is one odd cycle; for
even $n$ it consists of two cycles of length $n/2$, including two-edge
cycles when $n=4$. The binary relation facts give the exact alternatives:
\begin{itemize}
\item If $n$ is odd, completion fails exactly when all $s_i$ are equal.
\item If $n\equiv0\pmod4$, completion is always possible.
\item If $n\equiv2\pmod4$, each of the two variable cycles must contain
at least one constraint with $s_i\ne s_{i+1}$.
\end{itemize}
It is therefore enough to select $e_0$ on an odd cycle, $e_0,e_3$ on a
cycle of length two modulo four, and no edge otherwise. In the second
case $n\ge6$, the two selected edges are disjoint and control different
variable cycles. For a selected edge, the crossing row is
$L_i=2(s_i-s_{i+1})$, with four distinct variables. The gap condition
ensures the required unequal values.

\subsection{Completing a path}

Let an $S$ path have vertices $v_0,\ldots,v_{n-1}$ and $m=n-1$ edges,
and write $z_{-1}=z_m=0$. Then
\[
 R(v_i)/2=s_i+z_{i-1}+z_i,\qquad
 B(v_i)/2=d(v_i)-s_i-z_{i-1}-z_i.
\]
An isolated vertex needs no completion. For a two-vertex path, if the
ordinary endpoint degrees differ, at least one of the two pure colours
is legal: the two failures would require simultaneously
$s_0-s_1=0$ and $s_0-s_1=d(v_0)-d(v_1)$. If the degrees are equal,
select its only edge and require $s_0\ne s_1$.

For $n\ge3$, every nonendpoint has ordinary degree four. Edges between
such vertices again give~\eqref{eq:cyclebits}, producing two variable
paths, one on each parity of edge index. At a left endpoint of ordinary
degree four, put $\delta=s_0-s_1$. Its condition is
$\delta-z_1\ne0$: it fixes $z_1=1$ for $\delta=0$, fixes $z_1=0$ for
$\delta=1$, and otherwise imposes no condition. At a degree-three left
endpoint, its own bit $z_0$ must also be considered:
\begin{equation}\label{eq:threeendbits}
 z_0=1\Longrightarrow\delta-z_1\ne0,
 \qquad z_0=0\Longrightarrow\delta-z_1\ne-1.
\end{equation}
For $\delta=0$ this is equality of $z_0,z_1$; for $\delta=-1$ it
excludes $(0,0)$; for $\delta=1$ it excludes $(1,1)$; otherwise it is
unrestricted. In each case it is a full-projection relation joining the
two variable paths. The right endpoint has the same description after
reversing the indices. Thus:
\begin{center}\small
\begin{tabular}{ccl}
\toprule Endpoint degrees&$n$&Sufficient treatment\\\midrule
$3,4$&either parity&One variable path, at most one prescribed bit.\\
$4,4$&odd&Prescribed bits lie on different variable paths.\\
$4,4$&even&Require left $\delta\notin\{0,1\}$.\\
$3,3$&odd&A relation cycle with an even number $m-2$ of inequalities.\\
$3,3$&even&Require left $\delta\ne0$ to make a relation nonbijective.\\
\bottomrule
\end{tabular}
\end{center}
In the odd-order $3,3$ case, either a relation is nonbijective, or both
endpoint relations are equality and all $m-2$ others are inequality.
This includes the parallel endpoint relations when $n=3$.

Select the first internal edge only when $n$ is even and the endpoint
degrees are equal. For endpoints of degree three and $n\ge4$, the row
has four variables and is $L=2\delta$. For endpoints of degree four,
it has five variables and is $L=2\delta-1$, since the first vertex has
three crossing edges and the second has two. The bound $|L|\ge2$ in
the latter case actually gives $|L|\ge3$ and excludes $\delta=0,1$.
The two-vertex case was handled separately. Each path selects at most
one internal edge.

\begin{proof}[Completion of Theorem~\ref{thm:pathcyclepart}]
All selected $S$ edges form a matching. Combine their crossing rows
with the rows for the $T$ pairs. Each row has at least four distinct
variables. Each crossing variable appears at most once on each side,
because the selected edges and the $T$ pairs are matchings.
Lemma~\ref{lem:signedgap} supplies all crossing signs simultaneously.
The preceding completion arguments then colour each $S$ component.
The $T$ comparisons follow from~\eqref{eq:integerpairdifference}, and
all crossing edges are separated by parity. Local binary completion is
equivalently a two-literal clause system, solvable by strongly connected
components, so the whole construction is polynomial.
\end{proof}

\subsection{Why internal recolouring is a real extension}

Take $2t$ disjoint triangles as $S$ and a perfect matching $Q$ between
their vertices, with no pair inside a triangle, connecting the triangles
as units. Such a matching exists for every $t\ge1$: allocate the three
ports of each triangle according to a connected cubic loopless
multigraph on $2t$ units. For each pair $uv\in Q$ introduce $a,b\in T$
and edges $ua,ub,va,vb,ab$. This gives a connected nonbipartite graph
with $12t$ vertices and $21t$ edges, of degrees four on $S$ and three
on $T$. Theorem~\ref{thm:pathcyclepart} applies with all $ab$ mixed.

If all $S$ edges were blue while this mixed set stayed fixed, each
$K_{2,2}$ crossing block would require a vertex of $S$ with exactly one
red crossing edge. Otherwise both $S$ vertices send equal contributions
to the two ends of the mixed edge $ab$. The $3t$ blocks use disjoint
$S$ pairs and require at least $3t$ such vertices. In a wholly blue
triangle there can be at most one: two would both have blue degree six
on their common blue edge. There are only $2t$ triangles, a contradiction.
The new theorem restores legality by changing internal pure colours,
while preserving the mixed set and the original graph.

The partition requirement itself is not universal. In $K_{3,3}$ plus
one chord on one side, an induced matching covering $T$ can contain at
most two vertices. Choosing the chord leaves an $S$ vertex of internal
degree three; choosing a crossing edge leaves two ordinary degree-three
vertices of internal degree two. Taking $T=\varnothing$ also fails.
This six-vertex graph has a stored legal colouring; it is an applicability
obstruction, not a counterexample to the conjecture.

\section{Isolated vertices of the matching side (E3)}\label{sec:isolatepart}

Keep $d(v)\in\{3,4\}$ and the same conditions on $G[S]$, but now allow
$G[T]$ to consist of a matching and an internally isolated set $I$.
Let $C$ be the crossing graph. Let $F$ be the selected $S$ matching from
the path and cycle proof, and let $P$ be its set of endpoints.

\begin{theorem}[A protected parity join]\label{thm:protectedjoin}
Put $A=C[(S\setminus P)\cup I]$. If every component of $A$ contains
an even number of vertices of $I$, then $\lir(\doub G)\le2$.
The original $T$ matching remains mixed, and the only additional mixed
edges belong to $A$. The component condition is necessary and sufficient
for an additional edge set in $A$ with odd degrees on $I$ and even
degrees on $S\setminus P$.
\end{theorem}
\begin{proof}
The last assertion is the elementary parity-subgraph criterion. Necessity
follows by summing degrees in each component. For sufficiency, root a
spanning tree in each component and determine parent edges from leaves
upwards; the root equation follows from the even number of odd targets.
Let $N$ be the resulting edge set and make exactly $N$ and the original
$T$ matching mixed. Every $T$ colour degree is odd and every $S$ colour
degree is even, so crossing edges are automatically legal.

No $N$ edge touches a $T$ pair or a protected $S$ endpoint. Consequently
the rows for $T$ pairs and selected $S$ edges keep all their original
pure crossing variables, with no new constants. They still have at
least four terms and at most two occurrences per variable. Apply
Lemma~\ref{lem:signedgap} and then the preceding component completion.
At an unprotected vertex the even number of mixed crossing edges still
makes its crossing red contribution even; the local completion only
needs the protected row conditions. This proves all required comparisons.
\end{proof}

A component of $A$ with an odd number of $I$ vertices is a certificate
against this protected-join restriction. It does not certify that the
graph has no locally irregular two-colouring.

For example, choose odd $t\ge1$ and $r\ge0$. On $2t$ auxiliary vertices
take a connected bridgeless cubic loopless multigraph: three parallel
edges for $t=1$, or $C_{2t}$ with an antipodal matching for $t\ge3$.
Subdivide one edge by $2r$ new vertices. Use these auxiliary vertices as
$T$, replace each auxiliary edge by a new $S$ vertex joined to its ends,
and join the $3t+2r$ vertices of $S$ into one path. Pair the $2r$ new
$T$ vertices consecutively by added internal edges. This is a simple
graph with $5t+4r$ vertices and $9t+7r-1$ edges. Its $S$ path has odd
order, so $F=\varnothing$; the auxiliary graph with the designated edge
removed remains connected. Thus the main component of $A$ contains all
$2t$ isolated $T$ vertices and the theorem applies. For $r>0$, internally
paired and isolated $T$ vertices coexist with unbounded counts.

\subsection{An exact condition with an independent side}

Suppose now that $T$ is independent. Fix a crossing mixed set $N$ with
\begin{equation}\label{eq:fixedjoinparity}
 d_N(v)\equiv0\pmod2\ (v\in S),\qquad
 d_N(v)\equiv1\pmod2\ (v\in T),
\end{equation}
and require all other original edges to be pure. Call an $S$ component
critical if it is a cycle of length not divisible by four, or an
even-order path with both ordinary endpoint degrees three. All its
vertices have crossing degree two. Call it saturated when every one
of its vertices has $N$ degree two.

\begin{theorem}[Exact independent-side completion]\label{thm:exactindependent}
For the fixed $N$ satisfying~\eqref{eq:fixedjoinparity}, a pure completion
of all remaining edges exists if and only if no critical component is
saturated.
\end{theorem}
\begin{proof}
Let $c_v$ be crossing degree and $m_v=d_N(v)$. The crossing red
contribution divided by two can be selected independently at each $S$
vertex as any integer in
\begin{equation}\label{eq:crossinterval}
 s_v\in[m_v/2,c_v-m_v/2]\cap\mathbb Z.
\end{equation}
Independence holds because every crossing edge has only one $S$ end,
and $T$ has no internal comparison. At a nonisolated $S$ vertex the
possibilities are $(c,m)=(2,0),(2,2),(3,0),(3,2)$, giving respectively
$s\in\{0,1,2\},\{1\},\{0,1,2,3\},\{1,2\}$.

Saturation of a critical component forces $s_v=1$ throughout. On a
critical cycle, the step-two constraint graph has an odd cycle of
inequality relations. On an even-order path with degree-three ends,
the two endpoint equalities and the odd number of internal inequalities
are inconsistent. This proves necessity, including a two-vertex path.

For each unsaturated critical component choose one vertex with $m_v=0$,
set its $s$ to zero and all other $s$ values to one. On an odd cycle this
creates a nonzero neighbouring difference. On a cycle of length two
modulo four the two differences around the special vertex occur in
different variable cycles. On a path with degree-three ends the change
makes an endpoint or internal relation nonbijective (and directly
separates the two vertices when the path has order two). The earlier
relation argument gives a completion.

For a cycle of length zero modulo four, an odd-order path, or a path
with different endpoint degrees, take all $s_v=1$. An isolated vertex
can take any allowed value. The only remaining case is an even-order
path with both endpoint degrees four. Its endpoints can each take
values one and two. For order two choose these different values. For
order $n\ge4$, set the left endpoint and all internal values to one.
Set the right endpoint to one if $(n-4)/2$ is even, and to two otherwise.
The two unary constraints on one variable path are separated by exactly
$(n-4)/2$ inequalities, so this choice is consistent; the other variable
path has no conflicting unary constraints. Realise the chosen crossing
degrees using~\eqref{eq:crossinterval}, then complete the internal
components. All crossing edges satisfy parity separation.
\end{proof}

\subsection{A bridge obstruction and its repair}

The opposite parity configuration, developed more generally in
Section~\ref{sec:rootports}, gives the following useful special case.

\begin{corollary}[Even components with an independent opposite side]
\label{cor:evencomponents}
If $T$ is independent and every $S$ component is an even-order path or
an even cycle, under the degree-three/four assumptions above,
$\lir(\doub G)\le2$. There is no parity restriction on $|T|$.
\end{corollary}
\begin{proof}
Make an internal perfect matching of every $S$ component mixed, and
all other internal edges blue. Bipartition each component, and give its
two classes respectively zero and two wholly red crossing edges at each
vertex. There are at least two crossing edges at every vertex. Red
degrees are one and five; every internal red difference has absolute
value four, and its blue difference is nonzero because ordinary
degrees differ by at most one. All crossing edges are pure, joining
odd colour degrees in $S$ to even colour degrees in $T$.
\end{proof}

An infinite family shows why parity solvability alone does not guarantee
the first configuration. For $t\ge1$, subdivide one edge of
$C_{4t}$ plus its antipodal matching, obtaining a graph $Q_t$ of order
$4t+1$ with one degree-two vertex. Join three disjoint copies of $Q_t$
to a new centre through these vertices. The resulting cubic graph $H_t$
has three bridges. Let $T=V(H_t)$ and introduce an $S$ vertex on each
edge of $H_t$, giving the subdivision as the crossing graph. Denote the
three bridge vertices in $S$ by $w_1,w_2,w_3$. Add $w_1w_2$, make six
nonbridge $S$ vertices into a six-cycle, and pair all remaining $S$
vertices. There is an even number remaining, and there are enough
nonbridge vertices for every $t\ge1$.

This simple connected nonbipartite graph has $30t+10$ vertices and
$45t+18$ edges; exactly six vertices have degree four, and all others
degree three. Every $S$ component has even order, so
Corollary~\ref{cor:evencomponents} colours it. But if all $S$ internal
edges are pure, $S$ has even colour degrees, and $T$ odd colour degrees,
each half of a subdivided bridge is forced mixed. Indeed, its cut side
contains $4t+1$ odd targets in $T$, so summing the mixed degrees modulo
two forces the sole allowed crossing boundary edge into $N$. Thus the
crossing red contribution is two at both $w_1,w_2$. If their common
internal edge has weight $x\in\{0,2\}$, then
\[
 R(w_1)=R(w_2)=2+x,\qquad B(w_1)=B(w_2)=4-x.
\]
Its actual colour is illegal. The obstruction excludes that entire
parity configuration, not just a particular choice of $N$, while the
opposite configuration colours the same infinite family.

\section{Odd-component root ports at arbitrary degrees (E4)}\label{sec:rootports}

The following theorem removes the global degree restriction from
Corollary~\ref{cor:evencomponents} and permits odd paths, odd cycles and
internally isolated vertices. The needed port certificate is explicit.

\begin{theorem}[Root-port construction]\label{thm:rootports}
Let a finite simple graph have a supplied partition $S\sqcup T$ with
$T$ independent and $\Delta(G[S])\le2$. Suppose every nonisolated
vertex of $G[S]$ has at least two neighbours in $T$, and every $S$
internal edge $uv$ satisfies $|d(u)-d(v)|\le1$.

Select one root in each odd-order $S$ component and no root in an
even-order component. In an odd path, number vertices from an end and
select a root with even index; in a singleton, select its only vertex.
In an odd cycle, choose a root whose two neighbours are not both of
degree $d(r)+1$ and not both of degree $d(r)-1$. Let $\mathcal U$ be
the selected roots. Suppose a crossing edge set $N$ exists with:
\begin{itemize}
\item every edge of $N$ joins a root to $T$;
\item every $T$ vertex has even $N$ degree;
\item every internally nonisolated root has $N$ degree one or three;
\item every internally isolated root has positive odd $N$ degree.
\end{itemize}
Then $\lir(\doub G)\le2$, constructively. There is no bound on the
maximum degree of $G$ or on degrees in $T$; internally isolated
vertices of $S$ may include original leaves.
\end{theorem}
\begin{proof}
An eligible root exists on every odd cycle. Otherwise each vertex
would be a strict local minimum or maximum of the degree sequence,
with these two types alternating, which cannot close around an odd
cycle. Deleting the chosen root from an odd cycle leaves an even-order
path; deleting an even-index vertex from an odd path leaves two
even-order paths. Give each remaining path or even cycle an internal
perfect matching. Make all its matching edges and all edges of $N$
mixed; make every other $S$ internal edge blue.

Each nonroot $S$ vertex has one internal mixed edge and no mixed crossing
edge. Every root has an odd number of mixed crossing edges and no
internal mixed edge. Thus all $S$ colour degrees are odd, all $T$
colour degrees are even, and all crossing edges are automatically legal.

On each remaining even path or cycle, assign red degrees one and five
to the two bipartition classes. A nonroot has at least two pure
crossing edges, so these targets are realised by making zero or two
of them wholly red. Internal edges between nonroots have red difference
of absolute value four; their blue differences are nonzero because
ordinary degree differences are at most one.

Give each nonisolated root red degree three. If it has one port, use
one other crossing edge wholly red; if it has three ports, use no other
red crossing edge. The assumed two crossing neighbours ensure capacity
in the first case. For a root neighbour $u$, the internal edge $ru$ is
blue, and its blue-degree difference is respectively
\begin{equation}\label{eq:rootbluedifference}
 B(r)-B(u)=2(d(r)-d(u))-2\quad\text{if }R(u)=1,
 \qquad
 B(r)-B(u)=2(d(r)-d(u))+2\quad\text{if }R(u)=5.
\end{equation}
If $d(r)-d(u)=1$, label five is forced; if it is $-1$, label one is
forced; if it is zero, either works. On an odd path, the components
on the two sides of the root can be phased independently. On an odd
cycle, the two endpoints after root deletion must receive opposite
labels. The only inconsistent requests are that both must be one or
both must be five, exactly the two cases excluded by root eligibility.
Hence a compatible phase exists.

An internally isolated root needs no internal comparison: give all
its remaining crossing edges blue. Its colour degrees retain odd
parity, with no need for red degree three. Every crossing edge has
only one $S$ end, so all these choices are independent. The preceding
arguments cover every original edge and produce the colouring by
finitely many polynomial-time traversals.
\end{proof}

\begin{corollary}[An exact root parity test within capacity]\label{cor:rootparity}
Fix eligible roots $\mathcal U$, and let
$A=C[\mathcal U\cup T]$ in the crossing graph. A subgraph with odd
degrees on $\mathcal U$ and even degrees on $T$ exists if and only if
each component of $A$ contains an even number of roots. If every
internally nonisolated root has crossing degree at most four, this
condition guarantees the port set of Theorem~\ref{thm:rootports}.
\end{corollary}
\begin{proof}
Apply degree summation and spanning-tree parity recovery as in
Theorem~\ref{thm:protectedjoin}. An odd port count at a nonisolated
root of crossing degree at most four is necessarily one or three.
Internally isolated roots permit any positive odd count.
\end{proof}

For roots of larger crossing degree the parity condition is still exact
for odd/even edge selection, but it alone does not guarantee the
one-or-three capacity for nonisolated roots. A failed tree recovery with
five ports must not be called infeasibility. A supplied qualified port
set remains sufficient at arbitrary degrees.

Pairing roots with common neighbours in $T$ is a special case: give
each pair its two crossing ports, making all $T$ counts even. It is not
necessary. For odd $b\ge3$, take a vertex $s_0\in S$ and $b$ groups of
three $S$ vertices; in each group add a private triple of $T$ vertices,
join the two triples completely, and join $s_0$ to one chosen private
$T$ vertex. All $S$ components are singletons. The root adjacency graph
formed by common $T$ neighbours is $b$ copies of $K_4$ sharing $s_0$.
It has no perfect matching: deleting $s_0$ leaves $b$ odd components,
which cannot all be matched through a single vertex. Nevertheless the
root-port graph is connected and has $1+3b$ roots, an even number, so
Corollary~\ref{cor:rootparity} applies. Explicitly, in each group select
the edge from $s_0$, one other edge to that same private $T$ vertex,
and two edges from the other two roots to a common private $T$ vertex.
This also shows why internally isolated roots need arbitrary odd port
capacity. These particular examples are bipartite, whose ordinary
existence theorem is already known~\cite{GW23}; their purpose is to
separate two port criteria, not to claim a new bipartite result.

Finally, fixing the roots can still lose a legal colouring. Take
$S=\{0,1,2,3\}$, $T=\{4,5,6\}$ with ordered edges
\[
 04,15,24,26,35,36,23.
\]
The odd components give roots $0,1$; the permitted root-port graph
has two components each containing one root, so the root-only scheme
fails. Yet the weights $(1,1,1,2,1,0,0)$ give
\[
 R=(1,1,3,1,2,2,2),\qquad B=(1,1,3,5,2,2,2),
\]
which are legal edge by edge. Here the even component $2\! -\!3$
also receives crossing mixed ports. Thus the exact root parity test
does not solve general even-component transport or produce the required
partition for every graph.

\section{Evidence and limits of the integer partition results}\label{sec:partitionevidence}

The four new groups are not four independent solutions of the mother
conjecture. E1 contains a general signed tool and an unbounded-degree
paired-partition application. E2 enlarges the $S$ side to paths and
cycles within degrees three and four. E3 addresses some, but not all,
isolated vertices on the matching side and gives an exact fixed-interface
criterion when that side is independent. E4 handles odd components
under an arbitrary-degree root-port certificate; the even-component
existence corollary is retained as its simpler special case rather than
counted as an additional maximal theorem. Fixed mixed sets, port counts
and parity patterns are part of the statements, not interchangeable
notions of extension.

The implementation checkers test 647 signed systems, 15,757 local
path/cycle profiles against full binary enumeration, 52 independent-side
fixed-join domains, and 727 eligible fixed-root atlas domains. They also
check parameter families, the actual integer repairs and the explicitly
failed restricted strategies. Separate standard-library auditors,
without importing the constructors or solvers, rebuild the saved signed
sums, local truth tables, cuts and actual copy degrees. Across the four
auditors, the stored whole-graph witnesses account for 972 actual
copies: 100, 190, 258 and 424 in the respective archived source audits.
The release report supplies machine-readable per-audit counts
so that repeated or overlapping examples are not counted as distinct
new graph families. These are implementation and certificate audits;
the displayed proofs justify the unbounded parameter ranges.

The finite counterexamples to weakened signed-gap hypotheses are also
retained. With only three terms, antisymmetric incidence forms of $K_4$
would require every vertex to be a source or a sink to reach absolute
value at least two, forcing a bipartition. Allowing three occurrences
already fails for
\[
 L_1=y_0+y_1+y_2+y_3+y_4,\quad
 L_2=-y_0-y_1+y_2+y_3+y_4,\quad
 L_3=y_0+y_1-y_2-y_3+y_4.
\]
After a global sign change $L_1\ge3$, at most one variable is negative.
With none negative $L_2=1$; if the sole negative is in positions zero or
one then $L_3=-1$, otherwise $L_2=-1$. Arbitrary affine shifts fail too:
for antisymmetric incidence forms of $K_5$, $L_v+1=2d^+(v)-3$.
The condition $|L_v+1|\ge2$ allows only outdegrees $0,3,4$; at most one
vertex has outdegree zero, so their sum is at least twelve rather than
the ten edges of $K_5$. These are obstructions to proposed stronger
tools, not to locally irregular two-colourability of the doubled graphs.

No claim of exhaustive historical priority is made for these tools.
The known regular and bipartite cases remain attributed to their
original sources. General graphs still need a suitable partition,
compatible mixed-edge choices, and control of the remaining component
configurations. The unrestricted conjecture is not proved.
